\begin{pastexam}[label={px:nsysu-110-osds-02}]{中山資工所甲組 110 學年度 第 8 題}
（追蹤）

(15\%) In the hash method, suppose that a rehash function $h_i(x)$=(5$i$+$x$) mod 13 is applied
sequentially when a new element is inserted into the hash table, but a hash collision occurs. In
other words, $h_1$( ) is applied for the first time; $h_2$( ) is applied if a hash collision occurs;
$h_3$( ) is applied if a hash collision occurs again; and so on. What is the content of the hash
table after the elements 35, 22, 9, 24, 16, 19, 3, are inserted sequentially into an empty hash table
of size 13, indexed as 0 through 12 ?
\end{pastexam}

\begin{pastsoln}
答案：
\begin{center}
\begin{tabular}{*{13}{c}}
\toprule
0 & 1 & 2 & 3 & 4 & 5 & 6 & 7 & 8 & 9 & 10 & 11 & 12 \\
\midrule
19 & 35 & & 24 & & 3 & 22 & & 16 & & & 9 & \\
\bottomrule
\end{tabular}
\end{center}

$h_i(x) = (x + 5i) \bmod 13$，$i = 1, 2, 3, \ldots$：第一次看 $(x + 5) \bmod 13$，每碰撞一次就往後跳 $5$ 格（繞回
開頭）。 這是開放定址，探測序列的間隔固定是 $5$，本質上是間隔 $5$ 的線性探測（資結書第~12.3~節）。 $13$ 是質數，
$\gcd(5, 13) = 1$，所以 $13$ 次之內一定會看遍所有格子，只要表沒滿就找得到空格。 注意 $i$ 從 $1$ 開始，第一次
看的格子是 $(x + 5) \bmod 13$，不是 $x \bmod 13$。
\begin{center}
\begin{tabular}{rp{0.68\linewidth}c}
\toprule
$x$ & 依序探測的格子 $h_1, h_2, \ldots$ & 放在 \\
\midrule
35 & $40 \bmod 13 = 1$ & 1 \\
22 & $27 \bmod 13 = 1$（35 占了），$32 \bmod 13 = 6$ & 6 \\
9 & $14 \bmod 13 = 1$（占），$19 \bmod 13 = 6$（占），$24 \bmod 13 = 11$ & 11 \\
24 & $29 \bmod 13 = 3$ & 3 \\
16 & $21 \bmod 13 = 8$ & 8 \\
19 & $24 \bmod 13 = 11$（占），$29 \bmod 13 = 3$（占），$34 \bmod 13 = 8$（占），$39 \bmod 13 = 0$ & 0 \\
3 & $8$（占），$13 \bmod 13 = 0$（占），$18 \bmod 13 = 5$ & 5 \\
\bottomrule
\end{tabular}
\end{center}
共發生 $1 + 2 + 3 + 2 = 8$ 次碰撞。 第 $2, 4, 7, 9, 10, 12$ 格是空的。

如果把第一次誤算成 $h_0(x) = x \bmod 13$（$i$ 從 $0$ 開始），會得到完全不同的表（例如 35 放在第 $9$ 格），題目明確說
第一次用 $h_1$，要照題目。
\end{pastsoln}
